\section{Where the Weyl energy can go at the recognition endpoint}
\label{sec:endpoint-charge}

Write \(\lambda_1^\pm\le\lambda_2^\pm\le z_\pm\) for the eigenvalues
of the real Weyl blocks, set \(L=\max(z_+,z_-)\), and write
\(d\lambda^\pm=|W^\pm|^2_{\rm HS}\,dV\). Use the physical
volume-one Yamabe metric for \(D=\int|Z|^2dV\). Define the
conformal invariant
\[
 B=\int_M\left\{\tfrac13\sum_{\pm}(\lambda_1^\pm-\lambda_2^\pm)^2
                   +\min(z_+^2,z_-^2)\right\}dV .
\]
For any positive class on a nonstandard smooth homotopy sphere,
\begin{equation}\label{eq:endpoint-deficit}
                 D+3B < 8(e-T),\qquad e=E_+(g).
\end{equation}
Indeed \(|W^\pm|^2=3z_\pm^2/2+(\lambda_1^\pm-\lambda_2^\pm)^2/2\)
by trace freeness. Hence \(B=4e/3-\int L^2dV\).
Chen and Zhu's strict and equality classifications give
\(Y^2<36\int L^2dV\) on the indicated smooth type
\cite[Theorems 1.4 and 1.6]{ChenZhu2014Conformal}.
For the strict classification simple connectivity excludes every
nontrivial connected-sum factor by van Kampen. At equality the other
possibilities are the projective plane over \(\mathbb C\) and a
universal cover conformal to a product of constant-curvature surfaces.
The cover here is the compact simply connected manifold itself.
Both surface factors would be compact simply connected surfaces, so
the product has nonzero second homology. The projective plane also
has nonzero second homology. Substitution of \eqref{eq:budgets} gives
\eqref{eq:endpoint-deficit}.

For the spectral and curvature estimates suppose \(e_i\to T\)
on hypothetical nonstandard smooth types, which may vary with \(i\).
The subsequent disk selection concerns the fixed-type case \(I=T\).
These are steps under a contradiction hypothesis excluded in
Section~\ref{sec:endpoint-smoothing}. Put
\(y_*=8\sqrt3\,\pi\). Then \(D_i+B_i\to0\),
\(Y_i\to y_*\), and \(6\|L_i\|_{2,g_i}\to y_*\).
On the set where \(z_+\ge z_-\) the negative energy density is
the sum of \(3z_-^2/2\) and half its lower eigengap squared.
On the complementary set interchange the signs. The combined
energy of these two smaller blocks is at most \(3B_i/2\).
The sets in this estimate may vary with \(i\).

\begin{theorem}[Finite disks carrying both chiral energies]
\label{thm:endpoint-disks}
For every such sequence there are smooth positive-Ricci metrics
\(k_i\in[g_i]\), with a constant scale specified below, and a
subsequence having the following property. Given \(\epsilon,\zeta,r_*>0\),
there are finitely many fixed radii \(r_j<r_*\) and, at every
sufficiently large index, pairwise disjoint parametrized smooth
closed standard disks \(B_{ij}\subset B_{k_i}(x_{ij},r_j)\) satisfying
\begin{align}
 \int_{B_{k_i}(x_{ij},r_j)} |\Rm_{g_i}|_{\rm tens}^2\,dV_{g_i}
 &<\zeta\quad\text{for each }j,\label{eq:endpoint-curv}\\
 (\lambda_i^++\lambda_i^-)
 \big(M\setminus\textstyle\bigcup_j B_{ij}\big)&<\epsilon .
 \label{eq:endpoint-capture}
\end{align}
The centers tend to Euclidean regular points in one pointed limit of
the comparison metrics.
\end{theorem}

\begin{proof}
Let \(q_i\) be a volume-one modified Yamabe representative with
\(R_{q_i}-6L_{q_i}=c_i\) constant, where \(c_i\) is the infimum of
the modified Yamabe quotient. Its positive conformal factor is
\(C^{2,\alpha}\) for \(0<\alpha<1\), sufficient for the scalar
equation and test functions below
\cite[Lemma 2.1 and proof of Theorem 1.4]{ChenZhu2014Conformal}.
By their Theorem 1.7, \(c_i\le0\) on this smooth type. Put
\(l_i=\|L_{q_i}\|_{2,q_i}=\|L_{g_i}\|_{2,g_i}\). The
ordinary Yamabe inequality and Cauchy--Schwarz give
\[
 Y_i\le\int R_{q_i}\,dV_{q_i}
      =c_i+6\int L_{q_i}\,dV_{q_i}\le c_i+6l_i.
\]
Thus \(c_i\to0\) and \(0\le6(l_i-\int L_{q_i}\,dV_{q_i})
\le6l_i-Y_i\to0\). Since the volume is one,
\(\|L_{q_i}-l_i\|_2^2=2l_i(l_i-\int L_{q_i}\,dV_{q_i})\).
In particular \(\|R_{q_i}-Y_i\|_{2,q_i}\to0\).
The difference
\[
 \tfrac23(d\lambda_i^++d\lambda_i^-)-L_{q_i}^2\,dV_{q_i}
\]
is a positive measure of mass \(B_i\). The square of \(L_{q_i}\)
converges in \(L^1(q_i)\) to \((y_*/6)^2=4T/3\).
Consequently, for \(\mu_i=dV_{q_i}\),
\begin{equation}\label{eq:charge-volume-tv}
 \left\|\frac{\lambda_i^++\lambda_i^-}{2T}-\mu_i\right\|_{\rm TV}
 \longrightarrow0.
\end{equation}
Here the norm is the mass of the variation measure. Both Weyl
measures are those of the original physical class on its actual sets.
The identity uses the auxiliary volume \(dV_{q_i}\).

The same representative has a further curvature bound. The equation
gives \(R_{q_i}/6\) within \(o(1)\) of \(L_{q_i}\) uniformly and
hence \(R_{q_i}\to y_*\) in \(L^2(q_i)\). Gauss--Bonnet gives
\(\int|Z_{q_i}|^2=4e_i+\int R_{q_i}^2/12-32\pi^2\to0\).
It follows from the curvature decomposition and
\eqref{eq:charge-volume-tv} that
\[
 \big\|\,|\Rm_{q_i}|_{\rm tens}^2-64\pi^2\,\big\|_{L^1(q_i)}
 \longrightarrow0.
\]
Indeed the four-index Weyl contribution tends to \(8T\) in this
norm and the scalar contribution to \(y_*^2/6\). The identities
for \(q_i\) follow by \(C^2\) approximation of its positive
conformal factor on each compact source. If smooth representatives
are needed, choose the approximations separately at each index so
that the errors in these integrals and in volume tend to zero.
These curvature conclusions concern the auxiliary representative;
\eqref{eq:endpoint-curv} below concerns the physical Yamabe metric.

Choose a smooth positive-Ricci comparison in the class as in
Section~\ref{sec:physical-volume}. Normalize its scale by the least
radius carrying half of \(\mu_i\), set this radius to one, and use an
attaining center \(p_i\). This is legitimate since the smooth
comparison distance spheres have zero \(\mu_i\)-measure.
We give the changes to the proof of Theorem~\ref{thm:volume-balls}.
On a cutoff \(\chi\) of support of \(\mu_i\)-measure at most one half,
the Yamabe inequality implies
\[
 \big[(1-2^{-1/2})Y_i-\|R_{q_i}-Y_i\|_{2,q_i}\big]
 \|\chi\|_{4,q_i}^2 \le 6\|d\chi\|_{2,q_i}^2.
\]
The bracket has a uniform positive floor. The conformally invariant
fourth-power gradient estimate in \eqref{eq:physical-small-support}
therefore applies with this floor. The logarithmic cutoffs and
packing argument give \eqref{eq:physical-power} for \(\mu_i\)
and \eqref{eq:physical-noncollapse} for \(k_i\). In particular
the auxiliary volumes are tight and uniformly small on small
comparison balls, with uniform noncollapse of comparison volume
on each bounded window.

Write \(q_i=w_i^2 k_i\). The scalar equation gives
\(-\Delta_{k_i}w_i\le V_iw_i\), where
\(V_i=\tfrac16(R_{q_i})_+ w_i^2\), since \(R_{k_i}>0\).
On any comparison ball \(U\) one has
\[
 6\|V_i\|_{2,k_i;U}\le Y_i\sqrt{\mu_i(U)}
                         +\|R_{q_i}-Y_i\|_{2,q_i}.
\]
On each fixed bounded window choose the radius sufficiently small,
then the index sufficiently large, to absorb the potential term at
\(p=4\) in \eqref{eq:physical-moser}. The resulting local
\(L^8(k_i)\) bound on \(w_i\) uses only
\(\int w_i^4dV_{k_i}=1\). Packing covers the window by uniformly
many such half-balls. Thus the auxiliary volume densities \(w_i^4\)
are bounded in \(L^2(k_i)\) there. Against compactly supported
continuous tests in a common pointed realization, Cauchy--Schwarz
and noncollapsed comparison-volume convergence give
\(\mu=f\mathcal H^4\) with \(f\in L^2_{\rm loc}\) in the limit.
Tightness gives \(\mu(X)=1\).

The original measures \(\nu_i=|\Rm_{g_i}|_{\rm tens}^2dV_{g_i}\)
have bounded total mass by \eqref{eq:budgets}. Pass to a further
subsequence with a local weak limit for these measures. Euclidean
regular points outside the atoms of this limit have full
\(\mu\)-measure. At \(\mu\)-almost every such center the density ratio
\(\mu(B_r(x))/\mathcal H^4(B_r(x))\) tends to \(f(x)\in(0,\infty)\)
as in
\eqref{eq:physical-ordered-limits}. Apply the finite disjoint selection,
shrink and smooth embedding in the proof of
Theorem~\ref{thm:volume-balls} at these centers. For any
\(\delta'>0\) the disks capture more than \(1-\delta'\) of \(\mu_i\),
with the prescribed bound on the original \(\nu_i\) in the containing
balls. Take \(2T\delta'<\epsilon\) and transfer
\eqref{eq:charge-volume-tv} on those same source sets.
\end{proof}

At this endpoint the selection removes the missing complement estimate
for these disks. To use them in Theorem~\ref{thm:simultaneous} one
needs the actual bicollar charts and their aggregate cost. For
example, if the uncovered positive charge is less than \(\epsilon<T\),
a cost \(c_i<T-\epsilon\) yields \(e_i+c_i<2T\) for all large
\(i\), and a positive retained class of energy less than \(T\).
A bound on the curvature integral of a containing ball alone gives
no first-order comparison on the parametrized boundary.
